% ============================================================================
% preamble.tex — 《PyTorch 面试速通》
% ============================================================================

% --- 中文支持 (XeLaTeX) ---
\usepackage{ctex}
\usepackage{fontspec}

% --- 页面布局 ---
\usepackage[
  a4paper,
  top=2.5cm, bottom=2.5cm,
  left=2.4cm, right=2.4cm
]{geometry}

% --- 数学 ---
\usepackage{amsmath, amssymb, amsthm}
\usepackage{mathtools}
\usepackage{bm}
\usepackage{enumitem}

% --- 图表 ---
\usepackage{graphicx}
\usepackage{float}
\usepackage{booktabs}
\usepackage{longtable}
\usepackage{tabularx}
\usepackage{array}

% --- 颜色 ---
\usepackage{xcolor}
\definecolor{codebg}{RGB}{248,248,248}
\definecolor{codeframe}{RGB}{220,220,220}
\definecolor{kwcolor}{RGB}{0,0,180}
\definecolor{strcolor}{RGB}{163,21,21}
\definecolor{commentcolor}{RGB}{0,128,0}
\definecolor{numbercolor}{RGB}{128,0,128}

% --- 代码高亮 (listings, 零外部依赖) ---
\usepackage{listings}
\lstdefinestyle{pythonstyle}{
  language=Python,
  backgroundcolor=\color{codebg},
  basicstyle=\ttfamily\small,
  keywordstyle=\color{kwcolor}\bfseries,
  stringstyle=\color{strcolor},
  commentstyle=\color{commentcolor}\itshape,
  numberstyle=\tiny\color{gray},
  numbers=left,
  numbersep=8pt,
  frame=single,
  rulecolor=\color{codeframe},
  breaklines=true,
  breakatwhitespace=false,
  showstringspaces=false,
  tabsize=4,
  columns=flexible,
  keepspaces=true,
  xleftmargin=14pt,
  framexleftmargin=14pt,
  morekeywords={self,with,as,assert,yield,None,True,False,nn,torch,F},
}
\lstset{style=pythonstyle}
% 纯文本/Shell 代码
\lstdefinestyle{plainstyle}{
  backgroundcolor=\color{codebg},
  basicstyle=\ttfamily\small,
  frame=single,
  rulecolor=\color{codeframe},
  breaklines=true,
  showstringspaces=false,
  columns=flexible,
  xleftmargin=14pt,
  framexleftmargin=14pt,
}
\newcommand{\code}[1]{\lstinline[basicstyle=\ttfamily]|#1|}

% --- 交叉引用与超链接 ---
\usepackage{hyperref}
\hypersetup{
  colorlinks=true,
  linkcolor=blue!60!black,
  citecolor=green!40!black,
  urlcolor=cyan!50!black,
}
\usepackage[nameinlink]{cleveref}

% --- 彩色框 (面试题 / 速记 / 坑) ---
\usepackage[most]{tcolorbox}
\tcbuselibrary{breakable, skins}

% 面试题框（蓝）
\newtcolorbox{interview}[1][]{
  enhanced, breakable, arc=3pt,
  before skip=10pt, after skip=10pt,
  colback=blue!4, colframe=blue!45!black,
  boxrule=0.6pt, left=8pt, right=8pt, top=6pt, bottom=6pt,
  fonttitle=\bfseries, title=面试题, #1
}
% 速记卡框（绿）
\newtcolorbox{quickcard}[1][]{
  enhanced, breakable, arc=3pt,
  before skip=10pt, after skip=10pt,
  colback=green!4, colframe=green!45!black,
  boxrule=0.6pt, left=8pt, right=8pt, top=6pt, bottom=6pt,
  fonttitle=\bfseries, title=面试速记卡, #1
}
% 坑/注意框（橙）
\newtcolorbox{trap}[1][]{
  enhanced, breakable, arc=3pt,
  before skip=10pt, after skip=10pt,
  colback=orange!5, colframe=orange!55!black,
  boxrule=0.6pt, left=8pt, right=8pt, top=6pt, bottom=6pt,
  fonttitle=\bfseries, title=常见坑, #1
}
% 一句话答案框（灰）
\newtcolorbox{answer}[1][]{
  enhanced, breakable, arc=3pt,
  before skip=8pt, after skip=8pt,
  colback=gray!6, colframe=gray!45,
  boxrule=0.6pt, left=8pt, right=8pt, top=6pt, bottom=6pt,
  fonttitle=\bfseries, title=参考答案, #1
}

% --- 常用数学命令 ---
\newcommand{\R}{\mathbb{R}}
\newcommand{\E}{\mathbb{E}}
\DeclareMathOperator{\softmax}{softmax}
\DeclareMathOperator*{\argmax}{argmax}
\DeclareMathOperator*{\argmin}{argmin}
